% Hafta 1 — Sır bellekten nasıl dışarı sızar?
% Derleme: py -3.12 tools/latex/derle.py docs/week-1/assets/h01-21-bellekten-sizma-yollari.tex
\documentclass[tikz,border=8pt]{standalone}
\usepackage{cen429-cizim}
\begin{document}
\begin{tikzpicture}[node distance=24mm and 12mm]
  \node[baslik] (b) at (0,0) {Sır bellekten nasıl dışarı sızar?};
  \node[altbaslik] (alt) at (b.south west) {değişken kapsam dışına çıkar, baytlar kalır};

  \node[uyarikutu=70mm, below=16mm of alt.south west, anchor=north, xshift=45mm, text width=62mm] (hub)
    {\kb{Parola bellekte}\\kullanıldı, ama silinmedi};

  \node[kotukutu=34mm, below=26mm of hub.south west, anchor=north west, text width=28mm] (s1) {\kb{Çökme dökümü}\\core dump dosyası};
  \node[kotukutu=34mm, right=of s1, text width=28mm] (s2) {\kb{Takas / uyku}\\swap, hibernation};
  \node[kotukutu=34mm, right=of s2, text width=28mm] (s3) {\kb{Hata ayıklayıcı}\\sürece bağlanır};
  \node[kotukutu=34mm, right=of s3, text width=28mm] (s4) {\kb{Bellek yeniden kullanımı}\\başka kod eski içeriği okur};

  \draw[ok, draw=kotu] (hub.south) -- (s1.north);
  \draw[ok, draw=kotu] (hub.south) -- (s2.north);
  \draw[ok, draw=kotu] (hub.south) -- (s3.north);
  \draw[ok, draw=kotu] (hub.south) -- (s4.north);

  \node[dip, text width=160mm, below=9mm of s1.south -| s2, anchor=north]
    {Çözüm: \kod{explicit\_bzero} / \kod{SecureZeroMemory} --- derleyicinin kaldıramayacağı silme.};
\end{tikzpicture}
\end{document}
